% Folder 149, batch 05 — archivists' pages 81–100.
% Pass: Opus 5.5 (claude-opus-5-5), 2026-10-03 — first pass, unchecked against the pages by a human.
%
% All twenty pages are mathematics and all are transcribed. His own
% pagination 41–50 runs on pages 82, 84, 86, 88, 90, 92, 94, 96, 98, 100
% (the odd archivists' pages carry no number of his).
% Notation fixed for this batch: his script E for the fundamental-group
% extensions is set \mathcal{E} (\mathcal{E}_K, \mathcal{E}_{U_K}, …); his
% "Gm" is \mathbb{G}_m; T is kept as he writes it (the inertia/Tate group);
% the tilde objects (\tilde S, \tilde K, \tilde O) are those of the
% (non-strict) henselisation at s.
\documentclass[11pt,a4paper]{article}
\input{../preamble/grothendieck}

\folder{149}
\batch{5}
\pages{81}{100}
\foldertitle{[Autour de La "Longue Marche" à travers la théorie de Galois, pages 1 à 67] : notes manuscrites (s.d.).}
\dating{s.d. — le groupe « Autour de La "Longue Marche" à travers la théorie de Galois » (141 à 149) est daté [à partir de 1978]-1983}
\watermark{Édition de démonstration}

\begin{document}

\page{81}
\note{la page continue un argument commencé avant ce lot : l'extension (13) et les classes $\xi_i$ y ont été introduites plus haut.}

Dans le cas « arithmétique », on voit donc que l'ext.~(13) est scindée ssi
$\xi_i = 0$, i.e. ssi on peut trouver, globalement sur $S$, une uniformisante
sur $X$ le long de $g_i(S)$ ; et cette condition (toujours satisfaite
loc. \add{(Zar.)} \struck{sur} $S$) est en tous cas suffisante (cas arith. ou
pas !). Quand $\xi_i = 0$, parmi les scindages, il y a ceux provenant du choix
d'une \struck{section} \add{base} de $J_i/J_i^2$ qui soit
\uncertain{privilégiée}. L'indétermination du choix d'une telle base est dans
$\mathbb{G}_m(S) = H^0(S, \mathbb{G}_m)$ — celle du choix d'une section de (13)
est dans
\[
  (20)\qquad H^1(S, T) \simeq \varprojlim_n H^1(S, \mu_n) .
\]
On a ici des suites exactes de Kummer
\[
  0 \to \mathbb{G}_m(S)/\mathbb{G}_m(S)^n \to H^1(S, \mu_n) \to {}_n\mathrm{Pic}(S) \to 0
\]
\note{le premier terme de cette suite est surchargé et en partie biffé ; la lecture $\mathbb{G}_m(S)/\mathbb{G}_m(S)^n$ est celle que la limite qui suit demande, elle n'est pas lue sur la page.}
d'où par passage à la limite
\[
  (21)\qquad 0 \to \mathbb{G}_m(S)^{\wedge} \to H^1(S, T) \to \varprojlim_n {}_n\mathrm{Pic}\, S .
\]
Dans le « cas arithmétique », le troisième terme est nul, on trouve donc
\struck{$\mathbb{G}_m$} $H^1(S, T) \simeq \mathbb{G}_m(S)^{\wedge}$.

\page{82}
\note{p.~41 de l'auteur.}

Si le genre est zéro, prenant une de ces sections de $T$ sur $S$ comme section
à l'infini, OPS (quitte à localiser) $U_S = \mathbb{E}^1_S \smallsetminus T'$,
donc $f$ s'identifie à une section de $\mathbb{E}^1_S$ sur $S'$, i.e.
de $\mathcal{O}_S$ sur $S'$, donc (comme $\operatorname{codim}(Z, S) \geq 2$,
$S'$ \uncertain{normal}) elle se prolonge en une section de $\mathbb{E}^1_S$.
\uncertain{Et on conclut} comme précédemment, OK.

\struck{Considérons} Considérons donc les diviseurs irréductibles $D_i$ sur $S$,
ou ce qui revient au même, les pts $x_i$ de $X$ de codim~$1$, i.e. tels que
$\mathcal{O}_{x_i}$ soit un anneau de valuation discrète).
Considérons son normalisé $\overline{\mathcal{O}}_{x_i}$ dans $\overline{K}$,
choisissons un idéal maximal $\tilde{x}_i$ de $\overline{\mathcal{O}}_x$
\struck{\ill{}} (ces idéaux correspondent \add{au-dessus de $x$} aux pts du
normalisé $\tilde{S}$ de $S$ dans $\overline{K}$) et considérons son
stabilisateur $N_x$ dans $\mathcal{E}_K$, qui opère donc dans
$k(\tilde{x}) = \overline{k(x)}$, et s'envoie en fait, on le sait, \emph{sur}
$\operatorname{Gal}(\overline{k(x)}/k(x))$.
\note{il écrit ici $X$ et non $S$ (« les pts $x_i$ de $X$ ») ; transcrit tel quel.}

\page{83}
Soit $I_{\tilde{x}}$ le sous-groupe \uncertain{noyau} de l'h.m. obtenu (le
stabilisateur « géométrique » de $\tilde{x}$), d'où une suite exacte
\[
  (7)\qquad 1 \to I_{\tilde{x}} \to N_{\tilde{x}} \to \operatorname{Gal}(\overline{k(x)}/k(x)) \to 1 ,
\]
et par Kummer on trouve can.
\[
  (8)\qquad I_{\tilde{x}} \simeq T_{\infty}\bigl(\overline{k(x)}^{*}\bigr) \simeq T_{\infty}\bigl(\overline{K}^{*}\bigr)
\]
\marginal{à cause de l'hyp. de car.~$0$ !}
On notera que si $x$ est le pt génér. du diviseur $D$, $k(x)$ est le corps des
fonctions de $D$, i.e. un corps de type fini sur $\mathbb{Q}$.

Il est immédiat (sans supposer que le corps de base pour $S$ soit $\mathbb{Q}$)
que le noyau de $\mathcal{E}_K \to \mathcal{E}_S$ est le sous-groupe invariant
engendré par les $I_{\tilde{x}}$. Donc l'hypothèse que
$\mathcal{E}_K \to \mathcal{E}_{U_S}$ s'annule sur le dit noyau, signifie aussi
qu'il s'annule sur chacun des $I_{\tilde{x}}$. Soit donc $U_{\mathcal{O}_x}$
induit par $U$ sur $\operatorname{Spec} \mathcal{O}_x$, \ill{}
\note{la fin de la page, deux lignes rapides, n'est pas lue avec assurance ; l'argument continue en p.~84 sur un autre sujet, une page manque peut-être entre les deux, ou il y a là une coupure de l'auteur.}

\page{84}
\note{p.~42 de l'auteur.}

d'ailleurs $\mathbb{G}_m(S)$ n'a pas d'él. infiniment divisibles
\struck{\ill{}}, donc
\[
  (22)\qquad \mathbb{G}_m(S) \longrightarrow \mathbb{G}_m(S)^{\wedge} \simeq H^1(S, T)
\]
\marginal{(cas « arithm. »)}
est injectif. Ainsi, quand $\xi_i = 0$ i.e. quand (13) admet des scindages
« géométriques » (et il suffit pour cela qu'il existe une section
\uncertain{(A const.)}), ceux-ci forment un torseur sous \struck{\ill{}}
$\mathbb{G}_m(S)^{\wedge}$, qui s'identifie à un sous-torseur du
$\bigl(H^1(S, T) = \mathbb{G}_m(S)^{\wedge}\bigr)$-torseur de tous les
scindages de (13). Pour que la « description profinie de la géom. alg. absolue
sur $\mathbb{Q}$ » soit complète, il y faudrait également caractériser (en
termes de \add{cette} description profinie) ce sous-ensemble remarquable.

Je voudrais enfin comprendre aussi comment une section d'extension du type~(1)
peut se « spécialiser » en une du type~(2), dans le cas de sections de courbes
relatives. Pour ceci, je reprends

\page{85}
la sempiternelle situation de $U = X \smallsetminus T$ sur $S$, avec une
$g_i = g_{i_0}$ section de $T$ sur $\uncertain{I}$, et je suppose donnée une
section rationnelle \struck{$f$} $f$ de $U$ sur $S$, qui en un point
particulier $s$ de \struck{\ill{}} codim~$1$ va coïncider avec \uncertain{une}
$g_i$ en tant que section de $X$
\[
  g_i(s) = f(s)
\]
de sorte que l'h.m. induit sur le groupe d'inertie $I_{s'} \subset \mathcal{E}_K$
par le composé
\[
  \mathcal{E}_K \longrightarrow \mathcal{E}_{U_K} \longrightarrow \mathcal{E}_U
\]
(qui envoie en fait $I_{s'}$ dans $\pi$, et \uncertain{même} dans les $L_i$, ou
l'un d'eux) soit non \struck{nulle} triviale. Considérons donc le groupe de
décomposition $N_{s'}$ (cf. suite exacte (7), dont on a changé les notations),
qui s'identifie comme le $\pi_1$ de la « rondelle trouée » correspondant au
hensélisé (pas strict) de l'anneau $\mathcal{O}_{S,s}$ — ou mieux, le $\pi_1$
du fibré en

\page{86}
\note{p.~43 de l'auteur.}

disques troués autour du germe de variété
$D = \bigl(\overline{\{s\}}\bigr)_{\mathrm{red}}$ en $s$ …, et l'extension
de $N_{s'}$ par $\pi$ contenue dans l'ext. $\mathcal{E}_{U_K}$ de
$\mathcal{E}_K$ par $\pi_1$. La section donnée (plutôt, construite via $f_K$)
de l'ext. $\mathcal{E}_{U_K}$, en définit une de la sous-extension.
\struck{Si on considérer d'autres pts que scindages de}
\marginal{\struck{\ill{}} (passage biffé de traits verticaux)}
L'interprétation géométrique de ces constructions profinies est la suivante ;
soit \struck{\ill{}} $\tilde{\mathcal{O}}_s$ le hensélisé (pas strict) de
$\mathcal{O}_{S,s} = \mathcal{O}_s$ \add{(trait hensélien)}, et
$\tilde{S} = \operatorname{Spec}(\tilde{\mathcal{O}}_s)$, $\tilde{K}$ son corps
des fonctions, on a $U_{\tilde{S}}$ sur $\tilde{S}$, $U_{\tilde{K}}$ sur
$\tilde{K}$, et \struck{l'extension}
\[
  (23)\qquad N_{s'} \ \text{(groupe de décomposition de $\mathcal{E}_{K,\overline{K}}$ en $s'$)} \simeq \mathcal{E}_{\tilde{K}}
\]
et la \struck{suite} \add{$\mathcal{E}$} extension de ce groupe par $\pi$
\uncertain{induite} que
\[
  (24)\qquad 1 \to \pi \to \mathcal{E}_{U_{\tilde{K}}} \to \mathcal{E}_{\tilde{K}} \to 1 ,
  \qquad \pi = \pi_1(U_{\overline{K}}), \quad \mathcal{E}_{\tilde{K}} = \operatorname{Gal}(\overline{K}/\tilde{K}),
\]
s'envoyant dans
\[
  (25)\qquad 1 \to \pi \to \mathcal{E}_{U} \to \mathcal{E}_{\tilde{\mathcal{O}}} \to 1 ,
  \qquad \mathcal{E}_{\tilde{\mathcal{O}}} \simeq \operatorname{Gal}(\overline{k}/k)
\]
($k = k(s)$ étant le corps résiduel).
\note{dans (25) il écrit $\mathcal{E}_U$ ; au diagramme (26) de la page suivante le même terme est $\mathcal{E}_{U_{\tilde{\mathcal{O}}}}$.}

\page{87}
Ceci se traduit dans un diagramme de suites exactes, numéroté (26) :

\begin{tikzcd}[column sep=small, row sep=small]
 & & 1 \arrow[d] & 1 \arrow[d] & \\
 & & I \arrow[r, "\simeq"] \arrow[d] & I_{s'} \simeq T \arrow[d] & \\
1 \arrow[r] & \pi \arrow[r] \arrow[d, no head] & \mathcal{E}_{U_{\tilde{K}}} \arrow[r] \arrow[d] & \mathcal{E}_{\tilde{K}} \arrow[r] \arrow[d] & 1 \\
1 \arrow[r] & \pi \arrow[r] & \mathcal{E}_{U_{\tilde{\mathcal{O}}}} \arrow[r] \arrow[d] & \mathcal{E}_{\tilde{\mathcal{O}}} \arrow[r] \arrow[d] & 1 \\
 & & 1 & 1 &
\end{tikzcd}

\note{à côté de $\mathcal{E}_{\tilde{K}}$ il écrit $\simeq \operatorname{Gal}(\overline{K}/\tilde{K})$, et à côté de $\mathcal{E}_{\tilde{\mathcal{O}}}$, $\simeq \operatorname{Gal}(\overline{k}/k)$. L'égalité verticale entre les deux $\pi$ est marquée d'un « $\Vert$ ».}

le noyau $I$ de $\mathcal{E}_{U_{\tilde{K}}} \to \mathcal{E}_{U_{\tilde{\mathcal{O}}}}$
s'envoyant isomorphiquement par $\mathcal{E}_{U_{\tilde{K}}} \to \mathcal{E}_{\tilde{K}}$
sur le sous-groupe d'inertie $I_{s'}$, lui-même can. isom. à $T$. La section de
$\mathcal{E}_{U_{\tilde{K}}}$ sur $\mathcal{E}_{\tilde{K}} \simeq N_{s'}$,
induite par la section de départ de $\mathcal{E}_{U_K}$ sur $\mathcal{E}_K$
(elle-même définie par $f_K$ pt $K$-rat. de $U_K$ sur $K$) n'est autre bien sûr
que celle définie par le pt $\tilde{K}$-rationnel correspondant $f_{\tilde{K}}$
de $U_{\tilde{K}}$. Son composé avec
$\mathcal{E}_{U_{\tilde{K}}} \to \mathcal{E}_{U_{\tilde{\mathcal{O}}}}$ envoie
$I \simeq T$ dans $\pi$, i.e. on a un homomorphisme essentiellement de la forme
$k_i \circ (\mu \cdot \mathrm{id}_T)$, où $k_i \colon T \to \pi$ est un hom. de
lacets, relatif \uncertain{au groupe de décomposition} \ill{}, et l'indice $i$.
Ainsi $\mathcal{E}_{\tilde{K}}$ s'envoie dans le normalisateur de
$L_i = k_i(T)$, dans $\mathcal{E}_{U_{\tilde{\mathcal{O}}}}$. D'ailleurs

\page{88}
\note{p.~44 de l'auteur.}

la dernière ligne dans (26) s'identifie aussi à
\[
  (27)\qquad 1 \to \pi \to \mathcal{E}_{U_s} \to \mathcal{E}_{k(s)} \to 1
\]
associée à $U_s$, fibre de $U/S$ en $s$, sur le corps résiduel $k(s)$.

\textbf{Résumons}. La situation de la courbe relative $U/S$, et d'un
$s \in S$ de codim~$1$, définit par hensélisation (pas stricte) en $s$ le
diagramme (26) \struck{(25)}, dans la dernière ligne est can. isom. à (27)
(par l'\add{hom.} \struck{\ill{}} sur les $\pi_1$ correspondant :

\begin{tikzcd}[column sep=small, row sep=small]
U_s \arrow[r] \arrow[d] & U_{\tilde{\mathcal{O}}} \arrow[d] \\
s \arrow[r] & \operatorname{Spec} \tilde{\mathcal{O}}
\end{tikzcd}

). Une section rationnelle \add{$f$}/ de $U$ sur $S$, plus génér.
$\tilde{f}$/ de $U_{\tilde{\mathcal{O}}}$ sur $\tilde{S}$, définit une
section \add{$\mathcal{E}_{\tilde{f}}$}/ de la ligne médiane de (26), ce qui
revient au même un relèvement
\[
  (28)\qquad \lambda \colon \mathcal{E}_{\tilde{K}} \longrightarrow \mathcal{E}_{U_{\tilde{\mathcal{O}}}} \;\simeq\; \mathcal{E}_{U_s}
\]
de $\mathcal{E}_{\tilde{K}} \to \mathcal{E}_{\tilde{\mathcal{O}}}$
($\simeq \mathcal{E}_{k(s)}$) — où $\mathcal{E}_{\tilde{K}}$ par ailleurs
s'identifie au groupe de décomposition de $\mathcal{O}_{S,s}$ dans
$\overline{K}$. Les conditions suivantes sont
\note{dans (28) l'isomorphisme $\mathcal{E}_{U_s} \simeq \mathcal{E}_{U_{\tilde{\mathcal{O}}}}$ est écrit verticalement, sous le terme d'arrivée.}

\page{89}
équivalentes :
\begin{enumerate}
  \item[a)] cet hom. est injectif.
  \item[b)] La restriction de cet hom. à $I_{s'} \simeq T$ est de la forme
  $k_i \circ (\mu \cdot \mathrm{id}_T)$, où $k_i$ est un hom. lacet
  $T \to \pi$, et $\mu \in \mathbb{N}^{*}$.
  \item[c)] $\lambda \mid I$ n'est pas l'hom. trivial.
  \item[d)] la section envisagée \add{$\mathcal{E}_{\tilde{f}}$ de
  $\mathcal{E}_{U_{\tilde{K}}}$ sur $\mathcal{E}_{\tilde{K}}$} ne
  \struck{\ill{}} \add{provient} pas d'une section de
  $\mathcal{E}_{U_{\tilde{\mathcal{O}}}}$ sur $\mathcal{E}_{\tilde{\mathcal{O}}}$.
  \item[e)] La section rationnelle \struck{régulière} $\tilde{f}$ de
  $U_{\tilde{K}}$ sur $\tilde{S}$ n'est pas régulière (i.e. la sect. rat. $f$
  de $U$ sur $S$ n'est pas définie en $s$).
  \item[f)] \struck{Si} \add{Désignons par} $\tilde{f}$ \struck{désignons
  aussi} la section de $X_{\tilde{\mathcal{O}}}$ sur $\tilde{S}$ définie par
  $\tilde{f}$ \struck{\ill{}} \add{\ill{}}, on a
  $\tilde{f}(s) \in I = X_s \smallsetminus U_s$.
\end{enumerate}
De plus, dans ce cas le $i$ dans b) n'est autre que $\tilde{f}(s)$.

Nous plaçant dans ce cas, on pourra donc identifier le groupe de décomposition
$N_{s'} \simeq \mathcal{E}_{\tilde{K}}$ (ext. de $\mathcal{E}_{k(s)}$ par $\pi$)
à un sous-groupe de $\mathcal{E}_{X_s}$ \uncertain{(qui est)} d'indice fini
$\mu$ dans $N(L_i)$. Ainsi, notre scindage \uncertain{de
$\mathcal{E}_{U_K}$ sur $\mathcal{E}_K$}, \uncertain{ou} de
$\mathcal{E}_{U_{\tilde{K}}}$ sur $\mathcal{E}_{\tilde{K}}$, qui ne provient
pas d'un scindage de
\note{dans la parenthèse de la ligne précédente, « ext. de $\mathcal{E}_{k(s)}$ par $\pi$ » : on attendrait « par $T$ » (cf. p.~90) ; transcrit tel quel.}

\page{90}
\note{p.~45 de l'auteur.}

$\mathcal{E}_U$ sur $\mathcal{E}_S$ ni même de $\mathcal{E}_{U_{\tilde{\mathcal{O}}}}$
sur $\mathcal{E}_{\tilde{\mathcal{O}}}$, définit néanmoins, \struck{sur} le
sous-groupe $\mathcal{E}_{\tilde{K}} = N_{s'}$ de $\mathcal{E}_K$ (qui est
extension d'un $\mathcal{E}_k$ par un facteur \uncertain{\ill{}} maximal
$T \simeq I_{s'}$, où $k$ est un corps de degré de transcendance $n - 1$,
lorsque $S$ est de type fini sur un corps $k_0$, disons $\mathbb{Q}$) presque un
relèvement — seul le facteur $I_{s'} \simeq T$ y obstrue encore : en tous cas
$N_{s'} \simeq \mathcal{E}_{\tilde{K}}$ est plongé dans un
$N_{\mathcal{E}_{X_s}}(L_i)$, et toute section de cette extension de
$\mathcal{E}_{k(s)}$ par $T$ donne donc une section de $\mathcal{E}_{X_s}$ sur
$\mathcal{E}_{k(s)}$ — \emph{dirais-je} au pire.

Une \add{autre} façon de dire ces choses est la suivante. Une section
rationnelle $f$ de $U/S$ sur $S$, de domaine de définition l'ouvert $S'$ de $S$,
peut s'interpréter en termes de la famille des pts $f(s) \in U_s$, où $s$
parcourant l'un des pts de $S'$

\page{91}
(ou si on préfère, disons, l'un des pts fermés (car $S$ est de type fini sur un
corps, tel $\mathbb{Q}$ disons)). En termes de sections d'extension, on peut dire
que la section $\mathcal{E}_{f_K}$ de $\mathcal{E}_{U_K}$ sur $\mathcal{E}_K$
donne \uncertain{lieu à} une famille de sections \uncertain{voisines} : une
famille de sections des $\mathcal{E}_{X_s}$ sur les $\mathcal{E}_{k(s)}$ pour
$s \in S'$. \struck{Plus} \struck{\ill{} précisément} \add{précisément}
\ill{} cette correspondance \ill{} \struck{\ill{}} sections de
$\mathcal{E}_{U_K}$ sur $\mathcal{E}_K$, et sections des $\mathcal{E}_{X_s}$ sur
les $\mathcal{E}_{k(s)}$, une généralisation de ce qui se passe dans le cas d'un
point de codim~$1$ : On interprète le normalisé $S'$ de $S$ dans $\overline{K}$,
la fibre en $s$ s'identifie à l'ens. des idéaux maximaux de la clôture normale
de $\mathcal{O}_{S,s}$ dans $\overline{K}$, et si $N_{s'}$ désigne le
\struck{normalisateur} \add{stabilisateur} dans $\mathcal{E}_K$ d'un $s'$ dans
$S'_s$, on a encore une structure d'extension
\[
  (29)\qquad 1 \to I_{s'} \to N_{s'} \to \mathcal{E}_{k(s)} \to 1
\]
\note{ici $S'$ désigne le normalisé de $S$ dans $\overline{K}$, alors qu'à la page précédente $S'$ était le domaine de définition de $f$ ; le double emploi est sur la page.}

\page{92}
\note{p.~46 de l'auteur.}

où $I_{s'}$ \struck{est le groupe} $\simeq \mathcal{E}_{\tilde{K}}$, $\tilde{K}$
étant le corps des fractions du hensélisé \struck{\ill{}} strict de
$\mathcal{O}_{S,s}$. ($N_{s'}$ \struck{est celui du complété} \add{s'identifie
de même à} $\mathcal{E}_{\tilde{K}}$, où $\tilde{K}$ correspond au hensélisé
pas strict). Ceci posé, on trouve comme précédemment, si le composé
$\mathcal{E}_K \to \mathcal{E}_{U_K} \to \mathcal{E}_U$ (\uncertain{lequel} il
envoie $I_{s'}$ dans $\pi$) s'annule sur $I_{s'}$ (ce qui signifie justement que
$s \in S'$ …) \uncertain{qu'on a} un hom. $\mathcal{E}_{k(s)} \to \mathcal{E}_U$,
ou si on préfère
$\mathcal{E}_{k(s)} \to \mathcal{E}_{U_{\tilde{\mathcal{O}}_s}} \simeq \mathcal{E}_{X_s}$,
qui peut s'interpréter comme une section de $\mathcal{E}_{X_s}$ sur
$\mathcal{E}_{k(s)}$.

Et que se passe-t-il pour $s \notin S'$ ? \struck{\ill{}} On sait que
$D = S \smallsetminus S'$ est purement de codim~$1$ (c'est l'image inverse du
diviseur $T$ sur $X$ par la section $f$ de $X$ sur $S$, définie
\uncertain{hors codim.~$\geq 2$} …) \ill{} en fait le $\mu = \mu_\alpha$ de tout
à l'heure, relatif à
\note{dans la première ligne, « $I_{s'} \simeq \mathcal{E}_{\tilde K}$ » avec le hensélisé strict, puis $N_{s'} \simeq \mathcal{E}_{\tilde K}$ avec le hensélisé non strict : la même lettre $\tilde K$ sert aux deux corps.}

\page{93}
une composante $D_\alpha = \overline{\{s_\alpha\}}$ de ce diviseur, n'est autre
que sa multiplicité en $D_\alpha$. Ceci rappelé, on se demande donc ce qui se
passe pour les $s \in \bigcup D_\alpha$, \struck{ou} \add{et} \struck{voit que}
on regardera ceux qui sont $\in D_\alpha$, \add{pour $\alpha$ fixé,} i.e. qui
spécialisent un $s_\alpha$ fixé. On peut dire que la situation définit une
extension $\mathcal{E}_{U_{s_\alpha}}$ de $\mathcal{E}_{k(s_\alpha)}$ par $\pi$
— \uncertain{disons} comme tout pour $S$ et $K$, d'une courbe relative \ill{}
$U_{D_\alpha}$ sur $D_\alpha$ — ou $U_{s_\alpha}$ sur un ouvert régulier
convenable de $D_\alpha$, si on préfère — en passant au pt générique. Cette
fois-ci, pour la fibre générique, on n'a plus de \struck{section} pt.
\uncertain{rationnel}, on a un pt. rat. sur $k(s_\alpha)$ de $X_{s_\alpha}$ —
i.e. un pt rat. $i_\alpha \in X_{s_\alpha} \smallsetminus U_{s_\alpha}$, i.e.
qui est à l'infini \ill{} maintenant rat. sur $k_\alpha$. Mais on dispose d'un
$L_i \subset \pi$ relatif à cet $i_\alpha$, justement, \struck{de façon}
ess. canonique,

\page{94}
\note{p.~47 de l'auteur.}

qui nous intéresse justement : regarder la classe de scindages de
$\mathcal{E}_{X_s}$ sur $\mathcal{E}_{k(s)}$ relatifs à cet $L_i$ — et notamment
ceux relatifs au choix d'une uniformisante de $\mathcal{O}_{S,s_\alpha}$, i.e.
d'une équation locale de $D_\alpha$ au voisinage de son pt gén. $s_\alpha$.
\struck{(spécialisation)} (qui interviennent par le biais \uncertain{qu'elles
définissent} $U_{s_\alpha}/$ …). Quand on a ces sections-là, on voit comment
elles se spécialisent, du moins c'est clair en les pts de $D_\alpha$ en lesquels
la section rationnelle envisagée du faisceau conormal
$\mathcal{N}^{\vee}_{D_\alpha/S}$ (si on suppose $D_\alpha$ régulier) reste
régulière, \uncertain{et} du $\mathcal{N}_{D_\alpha/S}$ inversible. On peut se
proposer d'étudier \struck{comment} ce qui se passe en les autres points —
i.e. la spécialisation des sections d'extensions du deuxième type, en faisant
une étude de groupes profinis \uncertain{analogue} à la précédente, pour cette
question. Mais cela me semble d'un intérêt \uncertain{mineur} — \ill{}

\page{95}
plus une question \ill{} des courbes relatives globales, mais pour des
« éléments de courbe » relatifs — savoir, pour le complété formel d'une courbe
relative $X/S$ le long d'une section $g_i$, \add{ou} le diviseur de cette
section distinguée dans la structure.

Je vais plutôt revenir à la situation où on part d'une section de première
espèce d'un $\mathcal{E}_{U_K}$ sur $\mathcal{E}_K$, associée à un pt rationnel
$f_K$ de $U_K$ sur $K$. On peut donc dire (indépendamment de la donnée du modèle
$S$ de $K$ sur lequel $U_K$ pourrait déjà se définir \struck{\ill{}} en schéma
relatif élémentaire) que
\[
  (29)\qquad \varphi = \mathcal{E}_{f_K} \colon \mathcal{E}_K \longrightarrow \mathcal{E}_{U_K}
\]
la dite section de $\mathcal{E}_{U_K}$ sur $\mathcal{E}_K$ définit, pour toute
section de $\mathcal{E}_K$ sur $\mathcal{E}_{\mathbb{Q}} = \Gamma_{\mathbb{Q}}$,
\uncertain{ou plutôt} sur un sous-groupe ouvert $\Gamma'$ de celui-ci, une
section
\note{le numéro (29) est déjà employé p.~91 ; transcrit tel quel.}

\page{96}
\note{p.~48 de l'auteur.}

correspondante de $\mathcal{E}_{U_K}$ sur $\Gamma'$ — d'où, en particulier, des
actions \uncertain{extérieures} de $\Gamma'$ sur $\pi$. (Même chose d'ailleurs
pour des sous-corps $k'$ quelconques de $K$, \add{\ill{}} qui \uncertain{font}
apparaître $\mathcal{E}_K$ comme étant sur $\mathcal{E}_{k'}$ : la section (29)
donne une famille de sections de $\mathcal{E}_{U_K}$ sur des $\mathcal{E}_k$
\add{ou des sous-groupes ouverts $\Gamma'$ de ceux-ci}, \ill{} des sections de
$\mathcal{E}_K$ sur des $\mathcal{E}_k$ — ou des sous-groupes
$\Gamma' \subset \mathcal{E}_k$.) J'ai envie d'avoir une vue d'ensemble sur cette
famille de sections, et de la mettre en relation avec la famille de sections des
$\mathcal{E}_{X_s}$ sur les $\mathcal{E}_{k(s)}$ définie par (29), \add{ou
définissant des} actions extérieures (pas seulement \uncertain{ext.}) des
$\mathcal{E}_{k(s)}$ sur $\pi$. Pour ces derniers, les $k(s)$ \uncertain{n'}
apparaissent pas comme des sous-corps de $K$, mais plutôt \add{\ill{}} comme des
$\mathcal{O}_s/\mathfrak{m}_s$ — des corps résiduels de \struck{\ill{}} anneaux
locaux. Cependant, si on suppose $S$ de dim \struck{finie} \add{relative~$1$} sur

\page{97}
\marginal{On aurait intérêt ici de prendre pour $S$ une courbe régulière sur $k$ \ill{}}
un corps $k$, tous les $k(s)$ (pour $s \neq \eta$, pt gén.) apparaissent comme des
extensions finies de $k$, et tout $s \in S$ définit une classe de conjugaison
\uncertain{$\Delta_{s'}$} \ill{} des $\mathcal{E}_K$ (\ill{} \add{inertie}
relative sur $\mathcal{E}_k$ …), et les $N_{s'}$ \ill{} les sous-groupes de
décomposition relatifs, et on peut regarder les sections de l'extension des
$N_{s'}$ sur des $\mathcal{E}_{k(s)} \subset \mathcal{E}_k$ \ill{} (sections
« de deuxième espèce ») et les comparer avec la section donnée. Dans le cas où
la section rat. $f$ de $U_K$ sur $S$ est définie en $s$, il résulte de ce qui
précède que la restriction \uncertain{partielle} obtenue de $\mathcal{E}_{U_K}$
\ill{} classe de section $\mathcal{E}_{k(s)} \to \mathcal{E}_K$ est \uncertain{elle
aussi} de première espèce, et \struck{\ill{}} \add{à} \struck{\ill{}}
\uncertain{quasi-p.} \add{près} (action sur $\pi$) est induite par celle qui est
définie par le pt $f(s) \in U_s(k(s))$. Dans
\note{cette page est d'écriture rapide ; plusieurs liaisons entre les formules ne sont pas lues.}

\page{98}
\note{p.~49 de l'auteur.}

le cas contraire où $f$ n'est pas définie en $s$, on trouve une
\struck{section} section de 2\textsuperscript{e} espèce, normalisant un $L_i$
\struck{\ill{}} dans $\pi$.

À vrai dire, on n'a regardé jusqu'à présent que ce qui se passe en des
« places » de $K/k$ (corps des fonctions relatif d'une variété) où la
\struck{famille} courbe $U_K/K$ a « bonne réduction » — c'était explicite, on
partait d'une $U$ sur $S$ — et je voudrais examiner maintenant ce qui se passe
dans le cas où on sort de cette hypothèse : examiner quels types d'actions de
$\mathcal{E}_{k(s)} \subset \mathcal{E}_k$ on obtient par la construction
précédente, en les $s \in S$ où on a « mauvaise réduction ». Un cas type serait
celui où $k = \mathbb{Q}$, $K$ corps des fonctions d'une courbe modulaire pour
courbes elliptiques \struck{(\ill{} divisible, \ill{})} ponctuées, et $s$ place
à l'infini de $K$, correspondant à une courbe elliptique

\page{99}
dégénérée. On peut passer encore à l'hensélisé (non strict) de
$\mathcal{O}_s$, i.e. partir de la situation d'un trait hensélien $S$, à corps
résiduel $k$ (ici $\mathbb{Q}$ ou une extension finie de $\mathbb{Q}$), d'une
courbe $U_\eta$ de type $(g, \nu)$ (ici $(1,1)$) sur son pt générique $\eta$, on
s'intéresse à l'extension $\mathcal{E}_{U_\eta}$ de $\mathcal{E}_\eta$ par
$\pi$, où ici $\mathcal{E}_\eta$ s'insère encore dans
\[
  (30)\qquad 1 \to T \to \mathcal{E}_\eta \to \mathcal{E}_k \to 1
\]
\struck{\ill{} ici \ill{}}
\[
  (31)\qquad 1 \to \pi \to \mathcal{E}_{U_\eta} \to \mathcal{E}_\eta \to 1
\]
Ici l'action extérieure de $T$ sur $\pi$ n'est \uncertain{évidemment} pas
triviale (pas de « bonne réduction ») — donc le groupe $\mathcal{E}_k$
n'opère pas par lui-même extérieurement sur $\pi$. Mais tout scindage de (30),
\marginal{\uncertain{mais} \ill{} $\subset \pi$ \ill{} \ill{} \ill{} !
\ill{} — note en oblique dans la marge inférieure gauche, non lue}

\page{100}
\note{p.~50 de l'auteur.}

définit une extension de $\mathcal{E}_k$ par $\pi$, donc une action extérieure
— est-elle « admissible » ?, i.e. définie par une courbe algébrique ? Sans doute
pas — quelle courbe serait-ce, \uncertain{i.e. ce serait} la courbe « réduite »
de type $(g, \nu)$ (inexistante !) ?
\marginal{mais si !? \uncertain{très} \ill{} « réduits » de $\pi$ \ill{}}
Dès lors, si on a une section de (30) \ill{} une de (31) — qui \uncertain{comprend}
\uncertain{une} \add{affection} \ill{} de $\mathcal{E}_k$ sur $\pi$ —
\uncertain{dans} une action, qu'attendre d'une telle action ?

\struck{Sans} \ill{} m'y attendre, \ill{} à tâter \ill{} pour la structure :
l'$\infty$ des variétés modulaires de Teichmüller !

S'il faut hasarder des conjectures dans cette situation scabreuse, ce
pourrait-être celle-ci :

\textbf{Conjecture de bonne réduction} : Les conditions suivantes sont
équivalentes
\begin{enumerate}
  \item[a)] $U_\eta$ a bonne réduction sur $S$
  \item[b)] L'action extérieure de $T$ sur $\pi$ est triviale
\end{enumerate}
\note{la liste des conditions se poursuit au-delà de ce lot.}

\end{document}
